\newpage
\doublespacing
\section{Introduction} 
Empirical research at the intersection of the sociology of culture and stratification has been primarily concerned  with exploring inequalities linked to social position in access to and engagement with symbolic goods  \citep{dimaggio00, vaneijck00, katzgerro04}.  The primary question in this (exponentially growing) literature has been and continues to be centered on whether taste and culture consumption patterns---as constitutive of the lifestyles that mark Weberian status groups---continue to be associated with social position is the central question in this (now expansive) literature \citep{tomlinson03, katzgerro04}.  This organizing concern can be appreciated whether the construct of ``social position'' is conceptualized in terms of educational attainment \citep{bourdieu_darbel91, bryson96, vaneijck99}, membership in broadly defined occupational groups \citep{bourdieu84, dimaggio_useem78a, dimaggio_useem78b, peterson_simkus93, tomlinson03}, aggregate measures of social class \citep{katzgerro_shavit98, katzgerro02, katzgerro06}, or constructed scalar indices of social status based on patterns of social intercourse across occupational groups \citep{chan_goldthorpe05, chan_goldthorpe07a, chan_goldthorpe07b}. 

The main conclusion that has been reached in this still vibrant line of research is that lifestyle practices continue to be strongly shaped---albeit in increasingly complex ways---by social position \citep{holt98, katzgerro99, tomlinson03, emmison03, warde_etal08}. This runs counter to the predictions of ``postmodernization'' and class ``destructuration'' theorists who foresaw a swift decline in (and ultimate eradication of) the linkage between social position and lifestyle in market-dominated ``consumer societies''\citep[e.g.][]{baudrillard06, kingston00, pakulski_waters96}. In direct contrast to this prediction, we can confidently assert that generalized markers of social status and education (rather than social class narrowly defined) provide the best lens to understand cultural stratification patterns \citep{alderson_etal07,chan_goldthorpe07a, chan_goldthorpe07b, katzgerro02, katzgerro04, sullivan_katzgerro07, silva06, bennett_etal99, emmison03, ollivier04}.  Confirming the basic findings of an earlier generation of research \citep{baumol_bowen66, dimaggio_useem78a, ganzeboom82}. 

The most important empirical generalization that can gleaned from research program is this: as education and generalized occupational status increases, all forms of cultural participation also increase.  In essence, what it means to be high-status in late modern societies is to be an equal opportunity connoisseur of diverse cultural offerings and aesthetic experiences across domains (e.g. literature versus film), forms of delivery (e.g. non-profit versus commercial systems) and external legitimation (e.g. ``prescribed'' or ``authorized'' culture versus ``pop culture'', ``folk culture'' or ``kitsch'') \citep{peterson92, peterson_kern96, ollivier08, warde_etal08, sullivan_katzgerro07, vaneijck01, dimaggio_mukhtar04}.

\subsection{What is good taste good for?}
Nevertheless, it can be argued that the strong focus of the cultural stratification literature on the question of \emph{whether} lifestyle practices continue to be unevenly segmented according to class and status markers has resulted in a relative lack of attention to an equally important question:  what is the \emph{role} (or ``function'') of aesthetic consumption in the everyday lives of members of different social strata? More precisely, given that the highly educated and those who belong to high-prestige occupational groups (e.g. specializing in highly professionalized work in the cultural and technical spheres) are more likely to engage the consumption of cultural goods and the arts than those who belong to working-class and so-called routine service occupational groups, to what \emph{uses} do these cosmopolitan ``omnivores'' \citep{ollivier08} put the culture that they so avidly consume? In essence now that we know that eclectic tastes is the form that ``good'' taste takes in late-modern societies, it is time to ask what is [this kind of] good taste good for'' \citep{erickson91}?

This gap in our knowledge is surprising in light of the fact that the question of the role played by engagement with symbolic goods was an continues to be central to the theoretical tradition from which most analysts of the culture-stratification linkage depart in the first place \citep{collins75, dimaggio87, douglas_isherwood96, fine77}.  In this respect, Fine's basic complaint in relation to research of popular culture in the early 1970s, can be repeated vis a vis research on cultural stratification today.  This research, generally ignores the \emph{behavioral implications} implications of culture consumption, especially as they pertain to social interaction \citep[453]{fine77}. An important analytic lesson that can be gleaned from these foundational sources is that if we cannot get a solid theoretical handle on the uses to which persons put the cultural knowledge acquired through aesthetic consumption experiences then it is difficult to analytically justify the status of these symbolic resources as ``cultural capital'' in a given social setting \citep{mohr_dimaggio95, dimaggio96, dimaggio00, carter03}.  This is the reason why in their influential paper \citet{lamont_lareau88} linked the notion of cultural capital to an explicit conceptual framework that highlighted the active deployment of a given repertoire of cultural tastes and lifestyle practices for the purposes of drawing symbolic boundaries.  \citet{dimaggio87, dimaggio00} was also adamant---building on the early work of Collins \citeyearpar{collins75}---in proposing that aesthetic consumption of cultural goods must be connected to routine exploitation of cultural knowledge in social interaction within everyday settings (see \citet{erickson96} and \citet{lizardo06a} for empirical applications and theoretical elaborations of this idea). From this perspective, one substantively important way in which the set of cultural aptitudes that manifest themselves as cultural \emph{consumption} become actualized as cultural \emph{capital} is through the linkage of cultural engagement to everyday processes of informal socialization and thus status group reproduction in the Weberian sense \citep{dimaggio_mohr85, dimaggio96}.
